\CNsection{}{\CNwarn{} IMPORTANT DISCLAIMER}

\begin{CNquote}

\textbf{This module discusses technologies at various stages of maturity.} Many concepts presented here are speculative, under early research, or in preliminary standardization study phases. The 6G standard does not yet exist — ITU IMT-2030 framework is under development, and 3GPP has not begun formal 6G specification work (expected \textasciitilde{}2025-2028). Information is based on published research, 3GPP study items, and industry roadmaps as of 2024-2025.

\end{CNquote}

\Needspace{17\baselineskip}

\CNsubsection{}{Maturity Classification Used in This Module}

{\def\LTcaptype{none} % do not increment counter
\Needspace{13\baselineskip}
\begin{longtable}[c]{@{}
  >{\raggedright\arraybackslash\hspace{0pt}}m{(\linewidth - 2\tabcolsep) * \real{0.4375}}
  >{\raggedright\arraybackslash\hspace{0pt}}m{(\linewidth - 2\tabcolsep) * \real{0.5625}}@{}}
\toprule\noalign{}
\rowcolor{CNink}\CNth{Label} & \CNth{Meaning} \\
\CNheadrulesep
\endhead
\bottomrule\noalign{}
\endlastfoot
\textbf{STANDARDIZED} & Specified in 3GPP Technical Specifications, deployable \\
\textbf{STUDY ITEM} & Formally studied in 3GPP (TR documents), may become work items \\
\textbf{CANDIDATE} & Strong industry consensus, likely to enter standardization \\
\textbf{RESEARCH} & Active academic/\allowbreak{}industry research, no formal standardization yet \\
\textbf{SPECULATIVE} & Early-stage concepts, feasibility not fully demonstrated \\
\end{longtable}
}

\Needspace{14\baselineskip}

\CNsection{}{Overview: The Road to 6G}

\CNchained\CNsubsection{}{IMT-2030 Vision (ITU-R)}

The ITU-R Working Party 5D has outlined the IMT-2030 framework with these usage scenarios:

\begin{itemize}
\tightlist
\item
  \textbf{Immersive Communication} (evolution of eMBB)
\item
  \textbf{Hyper Reliable and Low-Latency Communication} (evolution of URLLC)
\item
  \textbf{Massive Communication} (evolution of mMTC)
\item
  \textbf{Ubiquitous Connectivity} (new)
\item
  \textbf{Integrated AI and Communication} (new)
\item
  \textbf{Integrated Sensing and Communication} (new)
\end{itemize}

\Needspace{26\baselineskip}

\CNsubsection{}{5G Capabilities vs 6G Targets (IMT-2030)}

{\def\LTcaptype{none} % do not increment counter
\Needspace{22\baselineskip}
\begin{longtable}[c]{@{}cccc@{}}
\toprule\noalign{}
\rowcolor{CNink}\CNth{KPI} & \CNth{5G (IMT-2020)} & \CNth{6G Target (IMT-2030)} & \CNth{Improvement} \\
\CNheadrulesep
\endhead
\bottomrule\noalign{}
\endlastfoot
Peak data rate & 20~Gbps & 200+ Gbps & 10\ensuremath{\times} \\
User experienced rate & 100~Mbps & 1-10~Gbps & 10-100\ensuremath{\times} \\
Latency (user plane) & 1~ms & 0.1~ms & 10\ensuremath{\times} \\
Reliability & 99.999\% & 99.99999\% & 100\ensuremath{\times} \\
Connection density & 10\textsuperscript{6} devices/\allowbreak{}km\textsuperscript{2} & 10\textsuperscript{7}-10\textsuperscript{8} devices/\allowbreak{}km\textsuperscript{2} & 10-100\ensuremath{\times} \\
Mobility & 500 km/\allowbreak{}h & 1000+ km/\allowbreak{}h & 2\ensuremath{\times} \\
Spectrum efficiency & 1\ensuremath{\times} (baseline) & 3-5\ensuremath{\times} & 3-5\ensuremath{\times} \\
Energy efficiency & 1\ensuremath{\times} (baseline) & 10-100\ensuremath{\times} & 10-100\ensuremath{\times} \\
Positioning accuracy & \textasciitilde{}1~m (R16) & 1-10~cm & 10-100\ensuremath{\times} \\
AI integration & External/\allowbreak{}OTT & Native/\allowbreak{}intrinsic & Paradigm shift \\
Sensing capability & Not supported & Integrated & New capability \\
\end{longtable}
}

\Needspace{17\baselineskip}

\CNsection{1}{AI-Native Networks /\allowbreak{} AI-RAN}

\CNchained\CNsubsection{}{Status: \CNdot{CNamber} STUDY ITEM (Release 18)}

\CNchained\CNsubsection{}{What}

AI-Native Networks integrate artificial intelligence as a fundamental design principle rather than an add-on optimization layer. AI/\allowbreak{}ML algorithms are embedded directly into the radio access network (AI-RAN) at every protocol layer.

\CNsubsection{}{Why /\allowbreak{} Problem Addressed}

\begin{itemize}
\tightlist
\item
  Traditional model-based algorithms (MMSE, LMMSE) reach performance limits in complex channels
\item
  Manual network optimization cannot scale to ultra-dense heterogeneous deployments
\item
  5G NR parameter space is enormous (beamforming, scheduling, power control)
\item
  Need for self-optimizing, self-healing networks
\end{itemize}

\CNkeepfig{m26-ai-layers}{3}

\CNsubsection{}{AI Integration Across Protocol Layers}

\CNfigure{m26-ai-layers}{AI across the protocol layers}

\CNsubsection{}{3GPP Progress}

\begin{itemize}
\tightlist
\item
  \textbf{TR 38.843}: Study on AI/\allowbreak{}ML for NR air interface (Release 18)
\item
  \textbf{Use cases studied}:

  \begin{itemize}
  \tightlist
  \item
    CSI feedback enhancement (AI-based compression)
  \item
    Beam management (AI-based beam prediction)
  \item
    Positioning enhancement (AI/\allowbreak{}ML-assisted)
  \end{itemize}
\item
  \textbf{TR 38.844}: Study on AI/\allowbreak{}ML for NG-RAN (network level)
\item
  \textbf{Key finding}: Performance gains demonstrated but deployment complexity is high
\end{itemize}

\CNsubsection{}{Benefits}

\begin{itemize}
\tightlist
\item
  10-30\% throughput improvement in complex scenarios
\item
  Reduced CSI feedback overhead (up to 50\% compression)
\item
  Faster beam acquisition (reduced beam sweeping)
\item
  Self-optimizing network behavior
\end{itemize}

\CNsubsection{}{Challenges}

\begin{itemize}
\tightlist
\item
  \textbf{Training data}: Requires massive, representative datasets
\item
  \textbf{Generalization}: Models may not transfer across environments
\item
  \textbf{Computational cost}: Inference latency vs. real-time requirements
\item
  \textbf{Standardization}: How to specify AI models in a standard?
\item
  \textbf{Interoperability}: Models from different vendors must coexist
\item
  \textbf{Explainability}: Black-box decisions in safety-critical scenarios
\end{itemize}

\CNsubsection{}{Relation to 5G}

\begin{itemize}
\tightlist
\item
  5G O-RAN architecture provides RIC framework (xApps/\allowbreak{}rApps)
\item
  R18 studies lay groundwork; actual specifications expected R19-R20
\item
  Evolution from ``AI-assisted'' (5G) to ``AI-native'' (6G)
\end{itemize}

\CNsubsection{}{Open Questions}

\begin{itemize}
\tightlist
\item
  How to handle model updates over-the-air?
\item
  What is the right split between model-based and data-driven?
\item
  How to ensure robustness against adversarial attacks?
\item
  Standardize the model or just the interface?
\end{itemize}

\Needspace{17\baselineskip}

\CNsection{2}{Extremely Large MIMO (XL-MIMO)}

\CNchained\CNsubsection{}{Status: \CNdot{CNred} RESEARCH}

\CNchained\CNsubsection{}{What}

XL-MIMO scales antenna arrays to thousands or tens of thousands of elements, creating electrically very large apertures. Unlike conventional massive MIMO (64-256 elements), XL-MIMO operates partly in the \textbf{near-field} regime, fundamentally changing propagation characteristics.

\CNsubsection{}{Why /\allowbreak{} Problem Addressed}

\begin{itemize}
\tightlist
\item
  Massive MIMO gains saturate with conventional array sizes
\item
  Need for extreme spatial multiplexing (100+ simultaneous streams)
\item
  Higher frequencies require larger arrays for equivalent beamforming gain
\item
  Support for extremely high data rates in 6G
\end{itemize}

\CNsubsection{}{Key Physics: Near-Field vs Far-Field}

Far-field boundary (Fraunhofer distance):

\CNdisplay{d_F = \frac{2D^2}{\lambda}}

Example: a $D = 2$ m array at 30~GHz ($\lambda = 10$ mm) gives $d_F = \dfrac{2 \times 4}{0.01} = 800$ m, so \textbf{most users are in the NEAR FIELD of large arrays!}

\textbf{Near-field characteristics:}

\begin{itemize}
\tightlist
\item
  Spherical wavefronts (not planar)
\item
  Beam focusing (not just steering) — energy concentrated at a point in space
\item
  Distance-dependent beamforming
\item
  Spatial non-stationarity: not all array elements ``see'' the same channel
\end{itemize}

\CNkeepfig{m26-visibility}{3}

\CNsubsection{}{Spatial Non-Stationarity}

\CNfigure{m26-visibility}{Spatial non-stationarity: each user sees only part of the array}

\CNsubsection{}{Benefits}

\begin{itemize}
\tightlist
\item
  Beam focusing enables spatial division even for users at same angle
\item
  Near-field MIMO provides additional degrees of freedom
\item
  Extreme spectral efficiency through massive spatial multiplexing
\item
  Inherent interference suppression
\end{itemize}

\CNsubsection{}{Challenges}

\begin{itemize}
\tightlist
\item
  Channel estimation complexity: $O(N^2)$ pilot overhead without structure exploitation
\item
  Hardware: cost, power, interconnect for thousands of RF chains
\item
  Signal processing: centralized processing infeasible \ensuremath{\to} distributed architectures
\item
  Channel modeling: 3GPP models assume far-field; new models needed
\item
  Calibration: mutual coupling, array imperfections at scale
\end{itemize}

\CNsubsection{}{Relation to 5G}

\begin{itemize}
\tightlist
\item
  Extension of 5G massive MIMO (R15-R17: up to 256 elements)
\item
  5G codebook-based and Type-II CSI approaches won’t scale
\item
  Requires new channel models beyond 3GPP TR 38.901
\end{itemize}

\CNsubsection{}{Open Questions}

\begin{itemize}
\tightlist
\item
  Optimal array geometry (planar, cylindrical, conformal)?
\item
  How to exploit near-field properties in standard protocols?
\item
  Hybrid analog/\allowbreak{}digital architectures at this scale?
\item
  Energy-efficient implementations?
\end{itemize}

\Needspace{17\baselineskip}

\CNsection{3}{Cell-Free Massive MIMO}

\CNchained\CNsubsection{}{Status: \CNdot{CNred} RESEARCH (5G-Advanced precursors exist)}

\CNchained\CNsubsection{}{What}

Cell-Free Massive MIMO eliminates the concept of cells. A large number of distributed Access Points (APs), connected via fronthaul to Central Processing Units (CPUs), coherently serve all User Equipments (UEs) in the area without cell boundaries.

\CNsubsection{}{Why /\allowbreak{} Problem Addressed}

\begin{itemize}
\tightlist
\item
  Cell-edge problem: users at cell boundaries suffer from interference and low SINR
\item
  Uneven QoS distribution in cellular networks (10th-percentile vs median)
\item
  Inter-cell interference is a fundamental limitation of cellular architecture
\item
  Need for uniform, ubiquitous coverage
\end{itemize}

\CNkeepfig{m26-cell-free}{3}

\CNsubsection{}{Architecture}

\CNfigure{m26-cell-free}{Cell-free massive MIMO: all APs serve all UEs}

\CNsubsection{}{Key Principles}

\begin{enumerate}
\def\labelenumi{\arabic{enumi}.}
\tightlist
\item
  \textbf{All APs serve all UEs} (in principle; practical systems use AP selection)
\item
  \textbf{Coherent joint transmission/\allowbreak{}reception} across distributed APs
\item
  \textbf{No handover} — UE is always served by the ``best'' cluster of APs
\item
  \textbf{User-centric clustering} — each UE has its own virtual cell
\end{enumerate}

\CNsubsection{}{Benefits}

\begin{itemize}
\tightlist
\item
  95th-percentile throughput improvement: 5-10\ensuremath{\times} over cellular
\item
  Elimination of cell-edge effect
\item
  Macro-diversity: robust against blockage and shadowing
\item
  Natural interference management through coordination
\item
  Uniform coverage quality
\end{itemize}

\CNsubsection{}{Challenges}

\begin{itemize}
\tightlist
\item
  \textbf{Fronthaul requirements}: Massive data exchange between APs and CPU
\item
  \textbf{Synchronization}: All APs must be phase-synchronized for coherent transmission
\item
  \textbf{Scalability}: Computational complexity grows with number of APs \ensuremath{\times} UEs
\item
  \textbf{Channel estimation}: Pilot contamination across large areas
\item
  \textbf{Practical deployment}: Cost of dense AP infrastructure
\end{itemize}

\Needspace{15\baselineskip}

\CNsubsection{}{Scalable Implementations}

{\def\LTcaptype{none} % do not increment counter
\Needspace{11\baselineskip}
\begin{longtable}[c]{@{}cccc@{}}
\toprule\noalign{}
\rowcolor{CNink}\CNth{Level} & \CNth{Processing} & \CNth{Fronthaul} & \CNth{Performance} \\
\CNheadrulesep
\endhead
\bottomrule\noalign{}
\endlastfoot
Level 1 & Fully centralized MMSE & Very high & Maximum \\
Level 2 & Local MMSE + central combining & High & Near-optimal \\
Level 3 & Local MF + central LSFD & Moderate & Good \\
Level 4 & Local MF, no cooperation & Low & Baseline \\
\end{longtable}
}

\CNsubsection{}{Relation to 5G}

\begin{itemize}
\tightlist
\item
  5G CoMP (Coordinated MultiPoint) is a limited precursor
\item
  5G-Advanced (R18): Multi-TRP enhancements move toward cell-free concepts
\item
  C-RAN architecture provides infrastructure basis
\item
  O-RAN fronthaul interfaces (F1, eCPRI) enable distributed processing
\end{itemize}

\CNsubsection{}{Open Questions}

\begin{itemize}
\tightlist
\item
  Optimal AP density and placement?
\item
  How to handle mobility efficiently?
\item
  Power control across distributed APs?
\item
  Economic viability vs. traditional deployment?
\end{itemize}

\Needspace{17\baselineskip}

\CNsection{4}{Integrated Sensing and Communication (ISAC /\allowbreak{} JCS)}

\CNchained\CNsubsection{}{Status: \CNdot{CNamber} STUDY ITEM (Release 19)}

\CNchained\CNsubsection{}{What}

ISAC (also called Joint Communication and Sensing, JCS) uses a shared waveform, shared hardware, and shared spectrum to simultaneously perform wireless communication and radar-like sensing (detection, ranging, velocity estimation, imaging).

\CNsubsection{}{Why /\allowbreak{} Problem Addressed}

\begin{itemize}
\tightlist
\item
  Separate radar and communication systems waste spectrum and hardware
\item
  Autonomous vehicles, drones, industrial robots need both connectivity and environmental awareness
\item
  Spectrum scarcity drives dual-use of allocated bands
\item
  5G NR waveform (OFDM) is already suitable for sensing with modifications
\end{itemize}

\CNkeepfig{m26-isac}{3}

\CNsubsection{}{Sensing Capabilities}

\CNfigure{m26-isac}{Integrated sensing and communication with one waveform}

Sensing parameters extracted:

\begin{itemize}
\tightlist
\item
  \textbf{Range:} $R = c\,\tau/2$ (from the delay $\tau$)
\item
  \textbf{Velocity:} $v = \lambda f_d / 2$ (from the Doppler shift $f_d$)
\item
  \textbf{Angle:} $\theta$ from array processing (AoA/\allowbreak{}AoD)
\end{itemize}

\CNkeepfig{m26-ofdm-radar}{3}

\CNsubsection{}{OFDM Radar Processing}

\CNfigure{m26-ofdm-radar}{OFDM radar processing on the resource grid}

\CNsubsection{}{3GPP Progress}

\begin{itemize}
\tightlist
\item
  \textbf{R19 Study Item}: ``Study on Integrated Sensing and Communication'' (approved 2024)
\item
  Key study areas:

  \begin{itemize}
  \tightlist
  \item
    Channel model for sensing (bistatic, monostatic)
  \item
    Waveform design and resource allocation
  \item
    Sensing-assisted communication (e.g., beam prediction from sensing)
  \item
    Privacy and regulatory aspects of network-based sensing
  \end{itemize}
\item
  Use cases: gesture recognition, object tracking, intrusion detection, weather monitoring
\end{itemize}

\CNsubsection{}{Benefits}

\begin{itemize}
\tightlist
\item
  Spectrum efficiency: dual use of same resources
\item
  Hardware reuse: shared antenna arrays and RF chains
\item
  Synergy: sensing assists communication (predictive beamforming) and vice versa
\item
  New services: network as a sensor (NaaS — Network as a Sensor)
\end{itemize}

\CNsubsection{}{Challenges}

\begin{itemize}
\tightlist
\item
  \textbf{Waveform tradeoff}: Optimal for communication \ensuremath{\neq} optimal for sensing
\item
  \textbf{Self-interference}: Monostatic sensing requires full-duplex capability
\item
  \textbf{Clutter}: Distinguish targets from environment
\item
  \textbf{Privacy}: Network can sense people without their devices
\item
  \textbf{Regulation}: Sensing in licensed communication bands needs new frameworks
\item
  \textbf{Performance metrics}: How to jointly optimize communication rate and sensing accuracy?
\end{itemize}

\CNsubsection{}{Relation to 5G}

\begin{itemize}
\tightlist
\item
  5G NR positioning (R16/\allowbreak{}R17) is a limited precursor
\item
  5G OFDM waveform is sensing-capable with proper processing
\item
  R19 study prepares for R20 work items
\item
  5G-Advanced adds framework; 6G makes it native
\end{itemize}

\CNsubsection{}{Open Questions}

\begin{itemize}
\tightlist
\item
  Monostatic vs bistatic vs networked sensing — which architecture?
\item
  How to allocate resources between sensing and communication optimally?
\item
  What sensing performance is achievable with communication waveforms?
\item
  Regulatory framework for passive sensing of people?
\end{itemize}

\Needspace{17\baselineskip}

\CNsection{5}{Reconfigurable Intelligent Surfaces (RIS)}

\CNchained\CNsubsection{}{Status: \CNdot{CNamber} STUDY ITEM (Release 18-19)}

\CNchained\CNsubsection{}{What}

Reconfigurable Intelligent Surfaces (RIS) are planar structures composed of many sub-wavelength elements that can be electronically controlled to reflect (or refract) incoming electromagnetic waves in desired directions. They reshape the wireless propagation environment without generating new signals.

\CNsubsection{}{Why /\allowbreak{} Problem Addressed}

\begin{itemize}
\tightlist
\item
  mmWave/\allowbreak{}sub-THz signals suffer from severe blockage
\item
  Traditional relays require full RF chains (expensive, power-hungry)
\item
  NLoS coverage gaps in dense urban environments
\item
  Need for energy-efficient coverage extension
\end{itemize}

\CNkeepfig{m26-ris}{3}

\CNsubsection{}{Operating Principle}

\CNfigure{m26-ris}{A reconfigurable intelligent surface steers the beam to a blocked UE}

\Needspace{15\baselineskip}

\CNsubsection{}{Types of RIS}

{\def\LTcaptype{none} % do not increment counter
\Needspace{11\baselineskip}
\begin{longtable}[c]{@{}
  >{\raggedright\arraybackslash\hspace{0pt}}m{(\linewidth - 6\tabcolsep) * \real{0.1500}}
  >{\raggedright\arraybackslash\hspace{0pt}}m{(\linewidth - 6\tabcolsep) * \real{0.4000}}
  >{\raggedright\arraybackslash\hspace{0pt}}m{(\linewidth - 6\tabcolsep) * \real{0.1500}}
  >{\raggedright\arraybackslash\hspace{0pt}}m{(\linewidth - 6\tabcolsep) * \real{0.3000}}@{}}
\toprule\noalign{}
\rowcolor{CNink}\CNth{Type} & \CNth{Characteristics} & \CNth{Gain} & \CNth{Complexity} \\
\CNheadrulesep
\endhead
\bottomrule\noalign{}
\endlastfoot
Passive RIS & Only phase shift, no amplification & Moderate & Low \\
Active RIS & Phase shift + amplification & High & Medium \\
STAR-RIS & Simultaneous transmit and reflect & High & High \\
Hybrid RIS & Some active, some passive elements & Configurable & Medium \\
\end{longtable}
}

\CNsubsection{}{Path Loss Model}

For a passive RIS with $N$ elements:

\CNdisplay{P_{\mathrm{Rx}} \propto N^2 \times L_{\mathrm{BS}\rightarrow\mathrm{RIS}} \times L_{\mathrm{RIS}\rightarrow\mathrm{UE}}}

\begin{itemize}
\tightlist
\item
  “Product-distance” path loss (multiplicative, not additive in dB)
\item
  \ensuremath{\to} the RIS must be placed near the BS or near the UE for best performance
\item
  \ensuremath{\to} the SNR gain scales as $N^2$ (coherent combining)
\end{itemize}

\CNsubsection{}{3GPP Progress}

\begin{itemize}
\tightlist
\item
  \textbf{TR 38.857}: Study on NR network-controlled repeaters (R18) — includes smart repeater concepts
\item
  \textbf{R19}: Further study on RIS-assisted NR
\item
  Study areas: channel modeling with RIS, deployment scenarios, control signaling
\end{itemize}

\CNsubsection{}{Benefits}

\begin{itemize}
\tightlist
\item
  Nearly passive operation (only control circuit needs power)
\item
  Low cost per element (printed circuit technology)
\item
  Flexible deployment (walls, ceilings, lamp posts)
\item
  Coverage extension without backhaul
\item
  Green communication: minimal energy consumption
\end{itemize}

\CNsubsection{}{Challenges}

\begin{itemize}
\tightlist
\item
  \textbf{Channel estimation}: Must estimate BS\ensuremath{\to}RIS and RIS\ensuremath{\to}UE channels separately
\item
  \textbf{Control overhead}: Configuring hundreds/\allowbreak{}thousands of elements in real-time
\item
  \textbf{Narrowband response}: Phase shifts are typically frequency-flat \ensuremath{\to} wideband limitation
\item
  \textbf{Deployment}: Optimal placement, orientation, density
\item
  \textbf{Multiplicative path loss}: Double path loss can limit practical gains
\item
  \textbf{Hardware impairments}: Discrete phase shifts, element coupling
\end{itemize}

\CNsubsection{}{Relation to 5G}

\begin{itemize}
\tightlist
\item
  5G network-controlled repeaters (R18) are a precursor
\item
  5G mmWave deployments highlight the blockage problem RIS addresses
\item
  R18-R19 studies inform potential R20 specifications
\item
  May complement (not replace) small cells and relays
\end{itemize}

\CNsubsection{}{Open Questions}

\begin{itemize}
\tightlist
\item
  How many elements needed for meaningful gain in practice?
\item
  Who controls the RIS — operator, third party, autonomous?
\item
  How to handle mobility (real-time reconfiguration speed)?
\item
  Passive vs active: which wins in cost-benefit analysis?
\item
  Can RIS work effectively for broadband (wideband) signals?
\end{itemize}

\Needspace{17\baselineskip}

\CNsection{6}{Sub-THz /\allowbreak{} THz Communications}

\CNchained\CNsubsection{}{Status: \CNdot{CNred} RESEARCH}

\CNchained\CNsubsection{}{What}

Sub-THz (100-300~GHz) and THz (300~GHz - 3~THz) communications exploit vast unused spectrum for extreme data rates. Channels of 10+ GHz bandwidth become available, potentially enabling 100+ Gbps per link.

\CNsubsection{}{Why /\allowbreak{} Problem Addressed}

\begin{itemize}
\tightlist
\item
  mmWave bandwidth (400~MHz - 2~GHz) insufficient for 6G peak rate targets (200+ Gbps)
\item
  Spectrum below 100~GHz is increasingly congested
\item
  Short-range applications (data kiosks, intra-device communication) need extreme throughput
\item
  Backhaul/\allowbreak{}fronthaul for dense networks needs wireless alternatives to fiber
\end{itemize}

\CNkeepfig{m26-thz-bands}{3}

\CNsubsection{}{Frequency Bands of Interest}

\CNfigure{m26-thz-bands}{Sub-THz and THz bands: bandwidth, range and applications}

Key atmospheric absorption windows:

\begin{itemize}
\tightlist
\item
  130–175~GHz (moderate absorption)
\item
  200–320~GHz (multiple windows between absorption lines)
\item
  water-vapour absorption lines at 183 and 325~GHz
\end{itemize}

\CNsubsection{}{Link Budget Challenge}

Free-space path loss at 300~GHz versus 30~GHz (same distance, 100~m):

\CNdisplay{\mathrm{FSPL} = 20 \log_{10}\!\left(\frac{4\pi d}{\lambda}\right)}

\begin{itemize}
\tightlist
\item
  at 30~GHz: FSPL = 91.5~dB
\item
  at 300~GHz: FSPL = 111.5~dB (\textbf{+20 dB!})
\item
  molecular absorption adds a further 1–10 dB/\allowbreak{}100~m (frequency-dependent)
\end{itemize}

Compensation requires very high-gain antennas (pencil beams), large arrays (small $\lambda$ \ensuremath{\to} many elements in a small area) and short distances.

\CNsubsection{}{Key Technologies Enabling THz}

\begin{enumerate}
\def\labelenumi{\arabic{enumi}.}
\tightlist
\item
  \textbf{III-V semiconductors} (InP, GaN): Power amplifiers up to 300~GHz
\item
  \textbf{Silicon-based (CMOS/\allowbreak{}SiGe)}: Lower cost, lower power, up to \textasciitilde{}200~GHz
\item
  \textbf{Photonic approaches}: Optical-to-THz conversion
\item
  \textbf{Graphene-based devices}: Theoretical potential for THz electronics
\item
  \textbf{Massive arrays}: 1000+ elements possible due to small wavelength
\end{enumerate}

\CNsubsection{}{Benefits}

\begin{itemize}
\tightlist
\item
  Enormous available bandwidth (10-100+ GHz channels)
\item
  Very small wavelength enables massive arrays in compact form factor
\item
  Inherent security (highly directional, short range)
\item
  High-resolution sensing capability (sub-mm resolution)
\end{itemize}

\CNsubsection{}{Challenges}

\begin{itemize}
\tightlist
\item
  \textbf{Severe path loss}: Requires LoS and very short range
\item
  \textbf{Atmospheric absorption}: Frequency-selective fading from molecular resonances
\item
  \textbf{Hardware}: Limited output power, high noise figure at THz frequencies
\item
  \textbf{Phase noise}: Oscillator stability at hundreds of GHz
\item
  \textbf{Blockage}: Even a hand can block a THz link completely
\item
  \textbf{Cost}: Exotic semiconductor processes
\end{itemize}

\CNsubsection{}{Relation to 5G}

\begin{itemize}
\tightlist
\item
  5G FR2 (24-52~GHz) demonstrated viability of high-frequency mobile access
\item
  Lessons from 5G mmWave beam management applicable to sub-THz
\item
  WRC-23 identified bands above 100~GHz for study (not yet allocated for mobile)
\item
  6G may initially target 100-300~GHz (sub-THz) before moving higher
\end{itemize}

\CNsubsection{}{Open Questions}

\begin{itemize}
\tightlist
\item
  Is mobile access viable above 100~GHz, or only fixed/\allowbreak{}short-range?
\item
  What is the maximum practical range for useful communication?
\item
  How to handle beam tracking at THz beam widths (\textless{}1\ensuremath{^{\circ}})?
\item
  Can silicon-based solutions achieve sufficient power at these frequencies?
\end{itemize}

\Needspace{17\baselineskip}

\CNsection{7}{Non-Terrestrial Networks (NTN)}

\CNchained\CNsubsection{}{Status: \CNdot{CNgreen} STANDARDIZED (basic R17) /\allowbreak{} \CNdot{CNamber} STUDY (advanced features)}

\CNchained\CNsubsection{}{What}

Non-Terrestrial Networks integrate satellite and aerial platforms into the 3GPP network architecture. NR and IoT-NTN enable direct communication between standard UEs (including smartphones) and LEO/\allowbreak{}MEO/\allowbreak{}GEO satellites or High-Altitude Platform Systems (HAPS).

\CNsubsection{}{Why /\allowbreak{} Problem Addressed}

\begin{itemize}
\tightlist
\item
  \textasciitilde{}90\% of Earth’s surface lacks terrestrial mobile coverage
\item
  Maritime, aviation, and remote areas have no connectivity
\item
  Disaster resilience requires non-terrestrial backup
\item
  Global IoT connectivity (agriculture, logistics, environmental monitoring)
\end{itemize}

\CNkeepfig{m26-ntn-altitudes}{3}

\CNsubsection{}{NTN Platform Types}

\CNfigure{m26-ntn-altitudes}{Non-terrestrial network platforms by altitude}

\Needspace{20\baselineskip}

\CNsubsection{}{3GPP Standardization Progress}

{\def\LTcaptype{none} % do not increment counter
\Needspace{16\baselineskip}
\begin{longtable}[c]{@{}ccc@{}}
\toprule\noalign{}
\rowcolor{CNink}\CNth{Release} & \CNth{Feature} & \CNth{Status} \\
\CNheadrulesep
\endhead
\bottomrule\noalign{}
\endlastfoot
R15-R16 & Study items (TR 38.811, 38.821) & Completed \\
R17 & Basic NR-NTN (transparent payload, GEO/\allowbreak{}LEO) & \textbf{STANDARDIZED} \\
R17 & IoT-NTN (NB-IoT/\allowbreak{}eMTC over satellite) & \textbf{STANDARDIZED} \\
R18 & NTN enhancements (mobility, coverage) & \textbf{STANDARDIZED} \\
R18 & NR-NTN on FR2 (above 10~GHz) & Study \\
R19 & Regenerative payloads, store-and-forward & Study \\
R19 & Direct-to-cell (smartphone satellite) & Study/\allowbreak{}WI \\
\end{longtable}
}

\CNsubsection{}{Key Technical Challenges Solved in R17}

\begin{enumerate}
\def\labelenumi{\arabic{enumi}.}
\tightlist
\item
  \textbf{Timing advance}: Large propagation delay (up to 540~ms for GEO) \ensuremath{\to} pre-compensation using GNSS at UE
\item
  \textbf{Doppler shift}: LEO at 7.5 km/\allowbreak{}s \ensuremath{\to} up to \ensuremath{\pm}24 ppm at S-band \ensuremath{\to} pre-compensation
\item
  \textbf{Large cells}: Satellite beams cover 100s of km \ensuremath{\to} timing advance differential within cell
\item
  \textbf{HARQ}: Disabled or modified due to long RTT (feedback arrives too late)
\end{enumerate}

\CNkeepfig{m26-ntn-payload}{3}

\CNsubsection{}{Architecture Options}

\CNfigure{m26-ntn-payload}{Transparent versus regenerative satellite payload}

\CNsubsection{}{Benefits}

\begin{itemize}
\tightlist
\item
  Ubiquitous global coverage
\item
  Disaster resilience and emergency communication
\item
  Direct-to-cell: unmodified smartphones connect to satellites
\item
  IoT coverage in remote areas (agriculture, shipping, pipelines)
\item
  Network redundancy and multi-connectivity
\end{itemize}

\CNsubsection{}{Challenges}

\begin{itemize}
\tightlist
\item
  \textbf{Link budget}: Very long distances require large antenna gains
\item
  \textbf{Spectrum coordination}: Coexistence with terrestrial networks
\item
  \textbf{Handover}: LEO satellites move rapidly (4-8~min visibility window)
\item
  \textbf{Interference}: Satellite beams overlap with terrestrial cells
\item
  \textbf{Capacity}: Limited per-user throughput from satellite
\item
  \textbf{Constellation cost}: Thousands of satellites needed for global LEO coverage
\end{itemize}

\CNsubsection{}{Relation to 5G}

\begin{itemize}
\tightlist
\item
  NTN is PART of 5G-Advanced (R17-R18 specifications exist)
\item
  Direct-to-cell (smartphone-satellite) commercially launching (T-Mobile/\allowbreak{}SpaceX, AST SpaceMobile)
\item
  6G targets tighter integration: seamless terrestrial/\allowbreak{}NTN handover, multi-layer NTN
\end{itemize}

\CNsubsection{}{Open Questions}

\begin{itemize}
\tightlist
\item
  How to achieve seamless handover between LEO satellites?
\item
  Regenerative vs transparent payload — which wins economically?
\item
  How much capacity can satellite provide per user?
\item
  Integration with terrestrial for load balancing in dense areas?
\end{itemize}

\Needspace{17\baselineskip}

\CNsection{8}{Semantic Communications}

\CNchained\CNsubsection{}{Status: \CNdot{CNred} RESEARCH /\allowbreak{} \CNdot{CNviolet} SPECULATIVE}

\CNchained\CNsubsection{}{What}

Semantic Communications transmit the \textbf{meaning} (semantics) of information rather than exact bit sequences. Instead of ensuring every bit is correctly received, the system ensures the intended meaning or task outcome is preserved, potentially requiring far fewer transmitted bits.

\CNsubsection{}{Why /\allowbreak{} Problem Addressed}

\begin{itemize}
\tightlist
\item
  Shannon’s theory optimizes bit transmission, not information utility
\item
  Most transmitted data has redundancy that classical coding doesn’t exploit
\item
  Video call: 90\%+ of pixels don’t change between frames
\item
  Machine-to-machine communication doesn’t need human-readable formats
\item
  Approaching Shannon capacity limits \ensuremath{\to} need new paradigm for further gains
\end{itemize}

\CNkeepfig{m26-semantic}{3}

\CNsubsection{}{Conceptual Framework}

\CNfigure{m26-semantic}{Classical versus semantic communication}

\CNsubsection{}{Key Technologies}

\begin{enumerate}
\def\labelenumi{\arabic{enumi}.}
\tightlist
\item
  \textbf{Deep Joint Source-Channel Coding (DeepJSCC)}

  \begin{itemize}
  \tightlist
  \item
    Single neural network replaces separate source and channel coding
  \item
    Graceful degradation (no cliff effect)
  \item
    Adapts to channel quality continuously
  \end{itemize}
\item
  \textbf{Knowledge Graphs /\allowbreak{} Shared Knowledge Base}

  \begin{itemize}
  \tightlist
  \item
    Transmitter and receiver share background knowledge
  \item
    Only novel/\allowbreak{}unexpected information needs transmission
  \item
    Dramatic compression for known contexts
  \end{itemize}
\item
  \textbf{Task-Oriented Communication}

  \begin{itemize}
  \tightlist
  \item
    Optimize for downstream task (classification, control) not reconstruction
  \item
    Transmit only task-relevant features
  \item
    Example: Autonomous driving \ensuremath{\to} transmit object positions, not full video
  \end{itemize}
\end{enumerate}

\Needspace{15\baselineskip}

\CNsubsection{}{Potential Bandwidth Savings}

{\def\LTcaptype{none} % do not increment counter
\Needspace{11\baselineskip}
\begin{longtable}[c]{@{}
  >{\raggedright\arraybackslash\hspace{0pt}}m{(\linewidth - 6\tabcolsep) * \real{0.2889}}
  >{\raggedright\arraybackslash\hspace{0pt}}m{(\linewidth - 6\tabcolsep) * \real{0.2444}}
  >{\raggedright\arraybackslash\hspace{0pt}}m{(\linewidth - 6\tabcolsep) * \real{0.2222}}
  >{\raggedright\arraybackslash\hspace{0pt}}m{(\linewidth - 6\tabcolsep) * \real{0.2444}}@{}}
\toprule\noalign{}
\rowcolor{CNink}\CNth{Application} & \CNth{Classical} & \CNth{Semantic} & \CNth{Reduction} \\
\CNheadrulesep
\endhead
\bottomrule\noalign{}
\endlastfoot
Text (translation task) & Full sentence bits & Intent + entities & 10-100\ensuremath{\times} \\
Image (classification) & All pixels (JPEG) & Feature vector & 100-1000\ensuremath{\times} \\
Video (surveillance) & Full frames (H.265) & Events + changes & 10-50\ensuremath{\times} \\
Control (robotics) & Sensor data stream & Action commands & 50-500\ensuremath{\times} \\
\end{longtable}
}

\emph{Note: These are theoretical/\allowbreak{}experimental estimates, not proven at scale}

\CNsubsection{}{Benefits}

\begin{itemize}
\tightlist
\item
  Extreme compression beyond Shannon limits for specific tasks
\item
  Graceful degradation under poor channel conditions
\item
  Natural integration with AI at endpoints
\item
  Efficient machine-to-machine communication
\item
  Reduced latency (fewer bits to transmit)
\end{itemize}

\CNsubsection{}{Challenges}

\begin{itemize}
\tightlist
\item
  \textbf{No mathematical framework}: Shannon theory doesn’t cover semantics
\item
  \textbf{Shared knowledge assumption}: Tx and Rx must agree on semantic model
\item
  \textbf{Generalization}: Different tasks need different semantic encoders
\item
  \textbf{Security}: Semantic attacks (adversarial examples that preserve bits but corrupt meaning)
\item
  \textbf{Standardization}: How to specify semantic protocols?
\item
  \textbf{Backward compatibility}: Cannot coexist with classical systems easily
\item
  \textbf{Evaluation metrics}: What replaces BER/\allowbreak{}BLER for semantic systems?
\end{itemize}

\CNsubsection{}{Relation to 5G}

\begin{itemize}
\tightlist
\item
  5G has no semantic communication features
\item
  Conceptually orthogonal to current standards
\item
  May first appear in application layer (not radio standard)
\item
  Could eventually change fundamental PHY design in 6G+
\end{itemize}

\CNsubsection{}{Open Questions}

\begin{itemize}
\tightlist
\item
  Is a universal semantic language possible, or always task-specific?
\item
  How to handle unknown/\allowbreak{}novel semantic content?
\item
  Can semantic and classical communication coexist on same channel?
\item
  Who defines the semantic model — standards body, AI training?
\item
  Timeline: Is this 6G or 7G technology?
\end{itemize}

\Needspace{17\baselineskip}

\CNsection{9}{Digital Twins for Networks}

\CNchained\CNsubsection{}{Status: \CNdot{CNred} RESEARCH}

\CNchained\CNsubsection{}{What}

A Network Digital Twin is a real-time, high-fidelity virtual replica of a physical network that continuously synchronizes with the live network. It enables what-if analysis, predictive optimization, and autonomous network management by simulating changes before deploying them.

\CNsubsection{}{Why /\allowbreak{} Problem Addressed}

\begin{itemize}
\tightlist
\item
  Network optimization is reactive (fix after failure)
\item
  Testing changes on live networks is risky
\item
  5G/\allowbreak{}6G network complexity exceeds human ability to manage
\item
  Need for predictive maintenance and proactive optimization
\item
  Training AI/\allowbreak{}ML models requires safe environments
\end{itemize}

\CNkeepfig{m26-digital-twin}{3}

\CNsubsection{}{Architecture}

\CNfigure{m26-digital-twin}{Network digital twin in a closed loop}

\CNsubsection{}{Use Cases}

\begin{enumerate}
\def\labelenumi{\arabic{enumi}.}
\tightlist
\item
  \textbf{Predictive maintenance}: Detect degradation before failure
\item
  \textbf{Configuration optimization}: Test parameter changes in simulation
\item
  \textbf{Capacity planning}: Simulate future growth scenarios
\item
  \textbf{AI/\allowbreak{}ML training}: Generate synthetic training data safely
\item
  \textbf{Anomaly detection}: Compare twin prediction vs reality
\item
  \textbf{Disaster recovery}: Pre-plan failover configurations
\item
  \textbf{Network slicing}: Simulate SLA compliance before provisioning
\end{enumerate}

\CNsubsection{}{Benefits}

\begin{itemize}
\tightlist
\item
  Risk-free experimentation on virtual network
\item
  Continuous optimization without service disruption
\item
  Faster AI model development and validation
\item
  Reduced OPEX through automation and prediction
\item
  Holistic network visibility
\end{itemize}

\CNsubsection{}{Challenges}

\begin{itemize}
\tightlist
\item
  \textbf{Fidelity}: Twin must accurately represent physical network behavior
\item
  \textbf{Real-time synchronization}: Data collection and model update latency
\item
  \textbf{Scale}: Modeling millions of devices in real-time is computationally expensive
\item
  \textbf{Validation}: How to verify the twin matches reality?
\item
  \textbf{Data requirements}: Enormous telemetry data streams
\item
  \textbf{Standardization}: No agreed data model or interface standard yet
\end{itemize}

\CNsubsection{}{Relation to 5G}

\begin{itemize}
\tightlist
\item
  5G O-RAN provides data interfaces (E2, O1, O2) that feed twins
\item
  5G network automation (SON) is a simpler precursor
\item
  Some vendors offer ``digital twin'' products today (mostly offline simulation)
\item
  6G targets real-time, closed-loop digital twins
\end{itemize}

\CNsubsection{}{Open Questions}

\begin{itemize}
\tightlist
\item
  What level of fidelity is ``good enough'' for different use cases?
\item
  How to handle model uncertainty and validation?
\item
  Centralized vs distributed twin architecture?
\item
  Standardized data models and APIs?
\end{itemize}

\Needspace{17\baselineskip}

\CNsection{10}{Advanced Positioning (cm-Level)}

\CNchained\CNsubsection{}{Status: \CNdot{CNamber} STUDY ITEM (Release 18 positioning enhancements)}

\CNchained\CNsubsection{}{What}

6G targets centimeter-level (1-10~cm) positioning accuracy using joint communication and positioning techniques, advanced signal processing, and AI/\allowbreak{}ML assistance. This goes far beyond GPS (\textasciitilde{}1-5~m) and 5G R16 positioning (\textasciitilde{}1~m outdoor).

\CNsubsection{}{Why /\allowbreak{} Problem Addressed}

\begin{itemize}
\tightlist
\item
  Indoor positioning remains unsolved (GPS unavailable)
\item
  Industrial automation requires mm-cm precision
\item
  Autonomous systems (drones, robots, vehicles) need precise localization
\item
  AR/\allowbreak{}XR applications require 6DoF tracking at cm level
\item
  Asset tracking in warehouses, hospitals, factories
\end{itemize}

\Needspace{15\baselineskip}

\CNsubsection{}{Positioning Techniques Evolution}

{\def\LTcaptype{none} % do not increment counter
\Needspace{11\baselineskip}
\begin{longtable}[c]{@{}ccc@{}}
\toprule\noalign{}
\rowcolor{CNink}\CNth{Release} & \CNth{Accuracy} & \CNth{Techniques} \\
\CNheadrulesep
\endhead
\bottomrule\noalign{}
\endlastfoot
R16 (5G baseline) & \textasciitilde{}3–10~m & DL-TDoA, UL-TDoA, multi-RTT \\
R17 (enhancement) & \textasciitilde{}0.5–3~m & Improved DL-PRS, NR positioning \\
R18 (study) & \textasciitilde{}0.2–1~m & Carrier phase, sidelink, AI/\allowbreak{}ML \\
R19+ /\allowbreak{} 6G target & \textasciitilde{}0.01–0.1~m & cm-level, 6DoF \\
\end{longtable}
}

\Needspace{11\baselineskip}

\CNsubsection{}{Key Technologies for cm-Level Accuracy}

\CNchained\CNsubsubsection{Carrier-Phase Positioning}

\begin{itemize}
\tightlist
\item
  \textbf{Standard ranging} uses the signal envelope: resolution $\sim c/BW$; for $BW = 100$ MHz the resolution is about 3~m (before processing gain).
\item
  \textbf{Carrier-phase ranging} uses the carrier wavelength $\lambda$: at 3.5~GHz, $\lambda = 8.6$ cm (sub-cm potential); at 28~GHz, $\lambda = 1.07$ cm (mm potential).
\item
  \textbf{Challenge:} integer-ambiguity resolution. The measured phase is
\end{itemize}

\CNdisplay{\varphi = \frac{2\pi}{\lambda}\,d + 2\pi N + \text{noise}}

and $N$ must be resolved to obtain the absolute distance.

\CNsubsubsection{Multi-RTT with Large Bandwidth}

\begin{itemize}
\tightlist
\item
  Round-Trip-Time measurement with 400~MHz - 2~GHz bandwidth
\item
  Sub-nanosecond timing resolution
\item
  Multiple base stations for triangulation
\end{itemize}

\CNsubsubsection{AI/\allowbreak{}ML-Assisted Positioning}

\begin{itemize}
\tightlist
\item
  \textbf{Fingerprinting}: ML models map channel features to position
\item
  \textbf{NLOS identification}: AI classifies LoS vs NLoS paths
\item
  \textbf{Sensor fusion}: Combine radio with IMU, camera, LiDAR
\item
  \textbf{Channel charting}: Unsupervised learning of radio geometry
\end{itemize}

\CNsubsubsection{Joint Communication and Positioning}

\begin{itemize}
\tightlist
\item
  Same waveform serves dual purpose
\item
  PRS (Positioning Reference Signals) integrated with data transmission
\item
  ISAC synergy: sensing capabilities enhance positioning
\end{itemize}

\Needspace{20\baselineskip}

\CNsubsection{}{Accuracy vs Technique Comparison}

{\def\LTcaptype{none} % do not increment counter
\Needspace{16\baselineskip}
\begin{longtable}[c]{@{}
  >{\raggedright\arraybackslash\hspace{0pt}}m{(\linewidth - 6\tabcolsep) * \real{0.3333}}
  >{\raggedright\arraybackslash\hspace{0pt}}m{(\linewidth - 6\tabcolsep) * \real{0.3030}}
  >{\raggedright\arraybackslash\hspace{0pt}}m{(\linewidth - 6\tabcolsep) * \real{0.1818}}
  >{\raggedright\arraybackslash\hspace{0pt}}m{(\linewidth - 6\tabcolsep) * \real{0.1818}}@{}}
\toprule\noalign{}
\rowcolor{CNink}\CNth{Technique} & \CNth{Accuracy} & \CNth{Pros} & \CNth{Cons} \\
\CNheadrulesep
\endhead
\bottomrule\noalign{}
\endlastfoot
DL-TDoA (R16) & 3-10~m & Standard, UE-based & Clock sync needed \\
Multi-RTT (R16) & 1-5~m & No sync needed & Multiple gNBs \\
DL-AoD (R16) & 5-15~m & Single gNB possible & Antenna cal needed \\
Enhanced (R17) & 0.5-3~m & Improved PRS & Bandwidth limited \\
Carrier-phase (R18) & 1-10~cm & Very high accuracy & Ambiguity resolution \\
AI-assisted (R18+) & 0.3-3~m & Works in NLOS & Training data needed \\
6G target & 1-10~cm & Native, real-time & Full system redesign \\
\end{longtable}
}

\CNsubsection{}{3GPP Progress}

\begin{itemize}
\tightlist
\item
  \textbf{R16}: Basic NR positioning (TS 38.305, 38.455)
\item
  \textbf{R17}: Positioning enhancements (improved accuracy, reduced latency)
\item
  \textbf{R18}: Study on AI/\allowbreak{}ML for positioning, carrier-phase, sidelink positioning
\item
  \textbf{R18}: Integrity and reliability requirements for positioning
\item
  \textbf{R19}: Further accuracy improvements, ISAC-based positioning
\end{itemize}

\CNsubsection{}{Benefits}

\begin{itemize}
\tightlist
\item
  Ubiquitous cm-level positioning (indoor and outdoor)
\item
  No additional infrastructure (reuses communication network)
\item
  6DoF tracking (position + orientation) for XR/\allowbreak{}AR
\item
  Enables new applications: digital twins, autonomous navigation
\end{itemize}

\CNsubsection{}{Challenges}

\begin{itemize}
\tightlist
\item
  \textbf{Carrier-phase ambiguity}: Integer ambiguity resolution is hard in mobile scenarios
\item
  \textbf{Multipath}: Reflections corrupt ranging measurements
\item
  \textbf{NLOS}: Dominant in indoor environments
\item
  \textbf{Synchronization}: Network timing accuracy limits achievable positioning
\item
  \textbf{UE complexity}: Processing requirements for high-accuracy positioning
\item
  \textbf{Power consumption}: Continuous positioning drains mobile battery
\end{itemize}

\CNsubsection{}{Relation to 5G}

\begin{itemize}
\tightlist
\item
  5G NR positioning (R16-R17) provides the foundation
\item
  Each release incrementally improves accuracy
\item
  6G integrates positioning as a native service (not an add-on)
\item
  Carrier-phase and AI techniques studied in R18 for potential R19 specification
\end{itemize}

\CNsubsection{}{Open Questions}

\begin{itemize}
\tightlist
\item
  Can carrier-phase work reliably in mobile environments?
\item
  How to handle integer ambiguity in real-time with mobility?
\item
  What infrastructure density is needed for cm-level indoor coverage?
\item
  How to balance positioning accuracy vs communication performance?
\end{itemize}

\Needspace{28\baselineskip}

\CNsection{}{Technology Maturity Matrix}

{\def\LTcaptype{none} % do not increment counter
\Needspace{22\baselineskip}
\begin{longtable}[c]{@{}
  >{\raggedright\arraybackslash\hspace{0pt}}m{(\linewidth - 8\tabcolsep) * \real{0.0526}}
  >{\raggedright\arraybackslash\hspace{0pt}}m{(\linewidth - 8\tabcolsep) * \real{0.1930}}
  >{\raggedright\arraybackslash\hspace{0pt}}m{(\linewidth - 8\tabcolsep) * \real{0.1053}}
  >{\raggedright\arraybackslash\hspace{0pt}}m{(\linewidth - 8\tabcolsep) * \real{0.2281}}
  >{\raggedright\arraybackslash\hspace{0pt}}m{(\linewidth - 8\tabcolsep) * \real{0.4211}}@{}}
\toprule\noalign{}
\rowcolor{CNink}\CNth{\#} & \CNth{Technology} & \CNth{TRL*} & \CNth{3GPP Status} & \CNth{Standardization Timeline} \\
\CNheadrulesep
\endhead
\bottomrule\noalign{}
\endlastfoot
1 & AI-Native Networks /\allowbreak{} AI-RAN & 4-5 & Study Item (R18 TR 38.843) & R19-R20 Work Items \\
2 & XL-MIMO & 2-3 & Not yet studied & R21+ (6G) \\
3 & Cell-Free Massive MIMO & 3-4 & Not yet studied (CoMP precursor) & R20-R21 \\
4 & ISAC & 3-4 & Study Item (R19) & R20 Work Items \\
5 & RIS & 3-4 & Study Item (R18-19 TR 38.857) & R20 Work Items \\
6 & Sub-THz Communications & 2-3 & Not yet studied & R21+ (6G) \\
7 & NTN (basic) & 7-8 & \textbf{Standardized} (R17) & Deployed \\
7 & NTN (advanced) & 4-5 & Study/\allowbreak{}WI (R18-19) & R19-R20 \\
8 & Semantic Communications & 1-2 & Not studied & R22+ (post-6G?) \\
9 & Digital Twins for Networks & 3-4 & Not yet studied & R20-R21 \\
10 & Advanced Positioning (cm) & 4-5 & Study Item (R18) & R19-R20 \\
\end{longtable}
}

*TRL = Technology Readiness Level (1 = basic research, 9 = deployed)

\Needspace{24\baselineskip}

\CNsubsection{}{TRL Scale Reference}

{\def\LTcaptype{none} % do not increment counter
\Needspace{20\baselineskip}
\begin{longtable}[c]{@{}cc@{}}
\toprule\noalign{}
\rowcolor{CNink}\CNth{TRL} & \CNth{Meaning} \\
\CNheadrulesep
\endhead
\bottomrule\noalign{}
\endlastfoot
1 & Basic principles observed \\
2 & Technology concept formulated \\
3 & Experimental proof of concept \\
4 & Technology validated in the lab \\
5 & Technology validated in a relevant environment \\
6 & Technology demonstrated in a relevant environment \\
7 & System prototype demonstrated in an operational environment \\
8 & System complete and qualified \\
9 & System proven in an operational environment (deployed) \\
\end{longtable}
}

\CNkeepfig{m26-roadmap}{5}

\CNsection{}{Projected Timeline: 3GPP Releases and 6G Technologies}

\CNfigure{m26-roadmap}{Projected timeline of 3GPP releases and 6G technologies}

\Needspace{15\baselineskip}

\CNsubsection{}{What to Expect per Release}

{\def\LTcaptype{none} % do not increment counter
\Needspace{11\baselineskip}
\begin{longtable}[c]{@{}
  >{\raggedright\arraybackslash\hspace{0pt}}m{(\linewidth - 4\tabcolsep) * \real{0.2093}}
  >{\raggedright\arraybackslash\hspace{0pt}}m{(\linewidth - 4\tabcolsep) * \real{0.2326}}
  >{\raggedright\arraybackslash\hspace{0pt}}m{(\linewidth - 4\tabcolsep) * \real{0.5581}}@{}}
\toprule\noalign{}
\rowcolor{CNink}\CNth{Release} & \CNth{Timeline} & \CNth{Key 6G-Related Features} \\
\CNheadrulesep
\endhead
\bottomrule\noalign{}
\endlastfoot
\textbf{R19} (5G-Adv) & 2025-2026 & ISAC study, advanced NTN, positioning enhancements, AI/\allowbreak{}RAN WI \\
\textbf{R20} (early 6G?) & 2027-2028 & ISAC specification, RIS specification, cell-free studies, sub-THz studies \\
\textbf{R21} (6G) & 2029-2030 & XL-MIMO, sub-THz, semantic study, full AI-native, digital twin \\
\textbf{R22} (6G+) & 2031-2032 & Semantic communication, advanced THz, autonomous networks \\
\end{longtable}
}

\Needspace{28\baselineskip}

\CNsection{}{Summary: 5G vs 6G Paradigm Shifts}

{\def\LTcaptype{none} % do not increment counter
\Needspace{22\baselineskip}
\begin{longtable}[c]{@{}
  >{\raggedright\arraybackslash\hspace{0pt}}m{(\linewidth - 4\tabcolsep) * \real{0.5000}}
  >{\raggedright\arraybackslash\hspace{0pt}}m{(\linewidth - 4\tabcolsep) * \real{0.2500}}
  >{\raggedright\arraybackslash\hspace{0pt}}m{(\linewidth - 4\tabcolsep) * \real{0.2500}}@{}}
\toprule\noalign{}
\rowcolor{CNink}\CNth{Aspect} & \CNth{5G} & \CNth{6G} \\
\CNheadrulesep
\endhead
\bottomrule\noalign{}
\endlastfoot
AI role & Optimization add-on & Native design principle \\
Spectrum & Sub-6 + mmWave (FR1+FR2) & + Sub-THz (FR3?) \\
Architecture & Cell-based & Cell-free + NTN \\
Coverage & Terrestrial & Terrestrial + space + aerial \\
Sensing & Not supported & Integrated (ISAC) \\
Positioning & \textasciitilde{}1m accuracy & \textasciitilde{}1~cm accuracy \\
Environment control & Passive (adapt to channel) & Active (RIS shapes channel) \\
Communication goal & Transmit bits reliably & Transmit meaning efficiently \\
Network management & Human-configured & Autonomous (digital twin + AI) \\
Antenna scale & 64-256 elements & 1000-10000+ elements \\
Peak rate & 20~Gbps & 200+ Gbps \\
Latency & 1~ms & 0.1~ms \\
\end{longtable}
}

\CNboxsection{Key Takeaways}

\begin{CNbox}[unbreakable]{sum}{Key Takeaways}

\begin{enumerate}
\def\labelenumi{\arabic{enumi}.}
\item
  \textbf{Not all ``6G technologies'' are equal} — NTN is already standardized, while semantic communication remains speculative. Always check maturity level.
\item
  \textbf{6G is evolutionary AND revolutionary} — Some features (AI-RAN, ISAC, RIS) evolve from 5G-Advanced. Others (semantic, THz, cell-free) represent paradigm shifts.
\item
  \textbf{3GPP drives standardization} — Technologies must pass: Research \ensuremath{\to} Study Item \ensuremath{\to} Work Item \ensuremath{\to} Technical Specification. This process takes 3-5 years per technology.
\item
  \textbf{Timing: 6G specification expected \textasciitilde{}2028-2030}, with first deployments \textasciitilde{}2030-2032.
\item
  \textbf{IMT-2030 framework} sets targets but actual specifications may differ from initial visions (as happened with IMT-2020/\allowbreak{}5G).
\item
  \textbf{Economic viability} ultimately determines adoption — technically superior solutions may lose to ``good enough'' practical alternatives.
\end{enumerate}

\end{CNbox}

\Needspace{14\baselineskip}

\CNsection{}{References and Further Reading}

\CNchained\CNsubsection{}{3GPP Documents}

\begin{itemize}
\tightlist
\item
  TR 38.843: Study on AI/\allowbreak{}ML for NR air interface (R18)
\item
  TR 38.857: Study on NR network-controlled repeaters (R18)
\item
  TR 38.811: Study on NR to support NTN (R15)
\item
  TR 38.821: Solutions for NR to support NTN (R16)
\item
  TS 38.305: Stage 2 functional specification for positioning (R16+)
\end{itemize}

\CNsubsection{}{ITU-R}

\begin{itemize}
\tightlist
\item
  IMT-2030 Framework Recommendation (ITU-R M.2160)
\item
  Report ITU-R M.2516: Future technology trends for IMT systems toward 2030+
\end{itemize}

\CNsubsection{}{Key Research Papers}

\begin{itemize}
\tightlist
\item
  Björnson et al., ``Scalable Cell-Free Massive MIMO Systems,'' IEEE Trans. Commun., 2020
\item
  Di Renzo et al., ``Smart Radio Environments Empowered by RIS,'' JSAC, 2020
\item
  Rappaport et al., ``Wireless Communications and Applications Above 100~GHz,'' IEEE Access, 2019
\item
  Qin et al., ``Semantic Communications: Principles and Challenges,'' arXiv, 2021
\item
  Liu et al., ``Integrated Sensing and Communications,'' IEEE JSAC, 2022
\item
  Amiri et al., ``Extremely Large MIMO,'' IEEE Commun. Mag., 2024
\end{itemize}

\CNsubsection{}{Industry Roadmaps}

\begin{itemize}
\tightlist
\item
  Next G Alliance (ATIS): 6G Roadmap
\item
  Hexa-X /\allowbreak{} Hexa-X-II (EU 6G Flagship Project)
\item
  Samsung 6G Vision (2020, updated 2023)
\item
  Nokia Bell Labs 6G Vision
\item
  NGMN Alliance: 6G Drivers and Vision
\end{itemize}

\emph{Module 26 — Last updated: August 2026}\CNbr{}\emph{Note: This is a living document. Technology maturity assessments will change as research progresses and 3GPP work advances.}
